\documentclass[12pt,a4paper]{article}

% ── Paquetes ──────────────────────────────────────────────
\usepackage[utf8]{inputenc}
\usepackage[T1]{fontenc}
\usepackage[spanish]{babel}
\usepackage{amsmath, amssymb, amsthm}
\usepackage{graphicx}
\usepackage{booktabs}
\usepackage{geometry}
\usepackage{hyperref}
\usepackage{xcolor}
\usepackage{listings}
\usepackage{float}

\geometry{margin=2.5cm}

% ── Colores para código Python ────────────────────────────
\definecolor{codebg}{RGB}{245,245,245}
\definecolor{codestr}{RGB}{163,21,21}
\definecolor{codecomment}{RGB}{0,128,0}
\definecolor{codekw}{RGB}{0,0,180}

\lstset{
  language=Python,
  backgroundcolor=\color{codebg},
  basicstyle=\ttfamily\footnotesize,
  keywordstyle=\color{codekw}\bfseries,
  stringstyle=\color{codestr},
  commentstyle=\color{codecomment}\itshape,
  showstringspaces=false,
  breaklines=true,
  frame=single,
  rulecolor=\color{gray!50},
  numbers=left,
  numberstyle=\tiny\color{gray},
  numbersep=6pt,
  xleftmargin=12pt
}

% ── Ambientes matemáticos ─────────────────────────────────
\newtheorem{definicion}{Definición}
\newtheorem{propiedad}{Propiedad}

% ── Portada ───────────────────────────────────────────────
\title{
  \textbf{Modelos (S)ARIMA para Series de Tiempo Financieras} \\[0.5em]
  \large Forma General, Herramientas de Diagnóstico y Aplicaciones en Python
}
\author{Análisis de Series de Tiempo}
\date{\today}

% ─────────────────────────────────────────────────────────
\begin{document}
\maketitle
\tableofcontents
\newpage

% ─────────────────────────────────────────────────────────
\section{Introducción}
% ─────────────────────────────────────────────────────────

Los modelos \textbf{ARIMA} (AutoRegressive Integrated Moving Average) y su extensión
estacional \textbf{SARIMA} son herramientas fundamentales para el análisis y pronóstico
de series de tiempo. En finanzas se utilizan para modelar precios de activos, tasas de
interés, retornos bursátiles y cualquier proceso que evolucione en el tiempo.

Este documento presenta: (1) la forma general del modelo $(S)ARIMA(p,d,q)(P,D,Q)_s$,
(2) los criterios de diagnóstico ACF, PACF, AIC y BIC, y (3) tres ejemplos prácticos
con datos financieros simulados en Python.

% ─────────────────────────────────────────────────────────
\section{Forma General del Modelo (S)ARIMA}
% ─────────────────────────────────────────────────────────

\subsection{Modelo ARIMA$(p,d,q)$}

Sea $\{y_t\}$ una serie de tiempo. El modelo \textbf{ARIMA}$(p,d,q)$ aplica $d$
diferencias ordinarias para obtener una serie estacionaria $w_t = \Delta^d y_t$,
que satisface:

\begin{equation}
  \phi_p(B)\,w_t = \theta_q(B)\,\varepsilon_t,
  \qquad \varepsilon_t \sim \mathrm{WN}(0,\sigma^2)
  \label{eq:arima}
\end{equation}

donde $B$ es el operador de rezago ($B\,y_t = y_{t-1}$),

\begin{align}
  \phi_p(B) &= 1 - \phi_1 B - \phi_2 B^2 - \cdots - \phi_p B^p
              \quad \text{(polinomio AR)}, \\
  \theta_q(B) &= 1 + \theta_1 B + \theta_2 B^2 + \cdots + \theta_q B^q
              \quad \text{(polinomio MA)}.
\end{align}

\begin{table}[H]
\centering
\caption{Parámetros del modelo ARIMA$(p,d,q)$}
\begin{tabular}{cll}
\toprule
\textbf{Parámetro} & \textbf{Significado} & \textbf{Identificación} \\
\midrule
$p$ & Orden autoregresivo   & PACF se corta en el rezago $p$ \\
$d$ & Número de diferencias & Prueba ADF / KPSS \\
$q$ & Orden de media móvil  & ACF se corta en el rezago $q$ \\
\bottomrule
\end{tabular}
\end{table}

\subsection{Extensión estacional: SARIMA$(p,d,q)(P,D,Q)_s$}

Cuando la serie presenta \textbf{estacionalidad de período $s$}, se añade una
componente estacional multiplicativa. La \textbf{forma general} es:

\begin{equation}
  \boxed{
    \Phi_P(B^s)\,\phi_p(B)\;\Delta_s^D\,\Delta^d\,y_t
    \;=\;
    \Theta_Q(B^s)\,\theta_q(B)\;\varepsilon_t
  }
  \label{eq:sarima}
\end{equation}

donde $\Delta_s = 1 - B^s$ es el operador de diferencia estacional y:

\begin{align}
  \Phi_P(B^s)   &= 1 - \Phi_1 B^s - \cdots - \Phi_P B^{sP}
                   \quad \text{(AR estacional)}, \\
  \Theta_Q(B^s) &= 1 + \Theta_1 B^s + \cdots + \Theta_Q B^{sQ}
                   \quad \text{(MA estacional)}.
\end{align}

\begin{table}[H]
\centering
\caption{Parámetros del modelo SARIMA$(p,d,q)(P,D,Q)_s$}
\begin{tabular}{cll}
\toprule
\textbf{Parámetro} & \textbf{Significado} & \textbf{Típico en finanzas} \\
\midrule
$p,d,q$   & Componente regular    & $d=1$ para precios \\
$P,D,Q$   & Componente estacional & $P,Q\in\{0,1\}$ \\
$s$       & Período estacional    & $s=12$ mensual; $s=4$ trimestral \\
\bottomrule
\end{tabular}
\end{table}

\subsection{Casos especiales}

\begin{itemize}
  \item \textbf{AR$(p)$}: $d=0,\;q=0$ — solo componente autoregresiva.
  \item \textbf{MA$(q)$}: $d=0,\;p=0$ — solo media móvil.
  \item \textbf{Caminata aleatoria}: ARIMA$(0,1,0)$ — modelo de referencia para precios.
  \item \textbf{ARIMA$(0,1,1)$}: equivale al suavizamiento exponencial simple.
\end{itemize}

% ─────────────────────────────────────────────────────────
\section{Herramientas de Diagnóstico e Identificación}
% ─────────────────────────────────────────────────────────

\subsection{Función de Autocorrelación (ACF)}

\begin{definicion}[ACF]
La \emph{función de autocorrelación} al rezago $k$ se define como:
\[
  \rho_k \;=\; \frac{\mathrm{Cov}(y_t,\,y_{t-k})}{\mathrm{Var}(y_t)}
           \;=\; \frac{\gamma_k}{\gamma_0}.
\]
\end{definicion}

Reglas de identificación:
\begin{itemize}
  \item \textbf{MA$(q)$}: la ACF se \emph{corta} bruscamente después del rezago $q$.
  \item \textbf{AR$(p)$}: la ACF decrece de forma exponencial o sinusoidal.
  \item Decaimiento muy lento $\Rightarrow$ serie no estacionaria ($d \geq 1$).
  \item Picos significativos en múltiplos de $s$ $\Rightarrow$ estacionalidad.
\end{itemize}

\subsection{Función de Autocorrelación Parcial (PACF)}

\begin{definicion}[PACF]
La \emph{función de autocorrelación parcial} al rezago $k$ mide la correlación entre
$y_t$ y $y_{t-k}$ eliminando el efecto lineal de los rezagos intermedios
$y_{t-1},\ldots,y_{t-k+1}$.
\end{definicion}

Reglas de identificación:
\begin{itemize}
  \item \textbf{AR$(p)$}: la PACF se \emph{corta} después del rezago $p$.
  \item \textbf{MA$(q)$}: la PACF decrece gradualmente.
\end{itemize}

\begin{table}[H]
\centering
\caption{Guía de identificación ACF--PACF (Box--Jenkins)}
\begin{tabular}{lll}
\toprule
\textbf{Modelo} & \textbf{ACF} & \textbf{PACF} \\
\midrule
AR$(p)$     & Decrece exponencial / sinusoidal  & Se corta en rezago $p$ \\
MA$(q)$     & Se corta en rezago $q$            & Decrece exponencial / sinusoidal \\
ARMA$(p,q)$ & Decrece después del rezago $q$    & Decrece después del rezago $p$ \\
\bottomrule
\end{tabular}
\end{table}

\subsection{Criterio de Información de Akaike (AIC)}

\begin{equation}
  \mathrm{AIC} = -2\,\hat{\ell} + 2k
\end{equation}

donde $\hat{\ell}$ es el log-verosimilitud maximizado y $k$ es el número de parámetros.
Se prefiere el modelo con \textbf{AIC menor}. Penaliza la complejidad para evitar sobreajuste.

\subsection{Criterio de Información Bayesiano (BIC)}

\begin{equation}
  \mathrm{BIC} = -2\,\hat{\ell} + k\,\ln(T)
\end{equation}

donde $T$ es el número de observaciones. El BIC penaliza la complejidad más fuertemente
que el AIC, favoreciendo modelos más parsimoniosos. En muestras grandes, BIC selecciona
el modelo ``verdadero'' con mayor probabilidad.

\subsection{Prueba de Dickey-Fuller Aumentada (ADF)}

La prueba ADF verifica:
\[
  H_0\colon \text{la serie tiene raíz unitaria (no es estacionaria)}.
\]
Si $p\text{-valor} < 0.05$ se rechaza $H_0$ y la serie es estacionaria.
Si no se rechaza, se aplica una diferencia ($d \leftarrow d+1$) y se repite.

% ─────────────────────────────────────────────────────────
\section{Proceso de Modelado Box--Jenkins}
% ─────────────────────────────────────────────────────────

El enfoque clásico sigue cuatro etapas:

\begin{enumerate}
  \item \textbf{Identificación}: verificar estacionariedad (ADF), transformar si es
        necesario ($\ln$, diferencias), inspeccionar ACF y PACF.
  \item \textbf{Estimación}: ajustar por máxima verosimilitud; comparar candidatos con
        AIC y BIC.
  \item \textbf{Validación}: analizar residuos $\{\hat{\varepsilon}_t\}$:
        ACF $\approx 0$ (Ljung-Box $p>0.05$), normalidad (Jarque-Bera),
        homocedasticidad.
  \item \textbf{Pronóstico}: generar predicciones con intervalos de confianza al 95\,\%.
\end{enumerate}

% ─────────────────────────────────────────────────────────
\section{Ejemplos en Python con Datos Financieros}
% ─────────────────────────────────────────────────────────

Los tres ejemplos usan datos simulados con propiedades estadísticas realistas
(tendencia, volatilidad y estacionalidad similares a las de los activos reales).
Para sustituir por datos reales basta reemplazar la sección de generación de datos
por una descarga con \texttt{yfinance}.

% ── Preámbulo Python común ────────────────────────────────
\begin{lstlisting}[caption={Importaciones y funciones auxiliares comunes a los tres ejemplos}]
import numpy as np
import pandas as pd
import matplotlib.pyplot as plt
from statsmodels.tsa.statespace.sarimax import SARIMAX
from statsmodels.graphics.tsaplots import plot_acf, plot_pacf
from statsmodels.tsa.stattools import adfuller
import warnings
warnings.filterwarnings('ignore')

np.random.seed(42)

def test_adf(serie, nombre):
    """Prueba Dickey-Fuller Aumentada."""
    res = adfuller(serie.dropna())
    print(f"ADF [{nombre}]: estadistico={res[0]:.4f}, p={res[1]:.4f}")
    if res[1] < 0.05:
        print("  -> Serie ESTACIONARIA (se rechaza H0)")
    else:
        print("  -> NO estacionaria; considerar diferenciacion")
    return res[1]

def graficar_acf_pacf(serie, titulo, lags=30):
    """Grafica la serie, ACF y PACF."""
    fig, axes = plt.subplots(1, 3, figsize=(15, 4))
    axes[0].plot(serie.values, linewidth=0.8)
    axes[0].set_title(f"{titulo} - Serie")
    plot_acf(serie.values,  lags=lags, ax=axes[1],
             title=f"ACF - {titulo}", zero=False)
    plot_pacf(serie.values, lags=lags, ax=axes[2],
             title=f"PACF - {titulo}", zero=False)
    plt.tight_layout()
    plt.savefig(f"fig_{titulo.replace(' ','_')}.png", dpi=120)
    plt.show()
\end{lstlisting}

% ─────────────────────────────────────────────────────────
\subsection{Ejemplo 1: Apple Inc.\ (AAPL) --- ARIMA$(2,1,2)$}
% ─────────────────────────────────────────────────────────

\textbf{Contexto.} Se modela el precio de cierre diario de Apple. La prueba ADF
confirma no estacionariedad en nivel ($p = 0.9013$) y estacionariedad tras la primera
diferencia ($p < 0.0001$), por lo que $d=1$. La inspección de ACF/PACF sugiere $p=2$,
$q=2$.

\begin{lstlisting}[caption={Ejemplo 1 --- ARIMA$(2,1,2)$ para AAPL}]
# ── Datos (simulados; reemplazar por yfinance para datos reales) ──
dates = pd.date_range("2020-01-01", "2024-12-31", freq='B')
n = len(dates)
trend    = np.linspace(300, 230, n)
seasonal = 15 * np.sin(2*np.pi*np.arange(n)/252)
noise    = np.random.randn(n).cumsum() * 2
aapl = pd.Series(trend + seasonal + noise, index=dates, name="AAPL")

# ── Estacionariedad ──────────────────────────────────────
test_adf(aapl,            "AAPL nivel")   # p = 0.9013 -> NO estacionaria
test_adf(aapl.diff(),     "AAPL 1a dif")  # p < 0.0001 -> estacionaria

# ── ACF y PACF ───────────────────────────────────────────
graficar_acf_pacf(aapl.diff().dropna(), "AAPL diff")

# ── Seleccion de orden por AIC ───────────────────────────
resultados = []
for p in range(4):
    for q in range(4):
        try:
            m = SARIMAX(aapl, order=(p,1,q), trend='n').fit(disp=False)
            resultados.append({'p':p,'q':q,'AIC':m.aic,'BIC':m.bic})
        except Exception:
            pass
df_aic = pd.DataFrame(resultados).sort_values('AIC').head(5)
print(df_aic.to_string(index=False))

# ── Ajuste ARIMA(2,1,2) ──────────────────────────────────
modelo = SARIMAX(aapl, order=(2,1,2), trend='n')
res    = modelo.fit(disp=False)
print(f"AIC = {res.aic:.2f}  |  BIC = {res.bic:.2f}")
print(res.params.round(4))

# ── Pronostico 30 dias ───────────────────────────────────
fc    = res.get_forecast(steps=30)
media = fc.predicted_mean
ic    = fc.conf_int(alpha=0.05)
print(f"Pronostico dia 30: {media.iloc[-1]:.2f} USD")
print(f"IC 95%: [{ic.iloc[-1,0]:.2f}, {ic.iloc[-1,1]:.2f}]")

# ── Grafica pronostico ───────────────────────────────────
plt.figure(figsize=(10,4))
plt.plot(aapl[-120:].values, label='Historico')
plt.plot(range(120, 150), media.values, color='red', label='Pronostico')
plt.fill_between(range(120, 150),
                 ic.iloc[:,0].values, ic.iloc[:,1].values,
                 alpha=0.2, color='red', label='IC 95%')
plt.title("AAPL - ARIMA(2,1,2): Pronostico 30 dias")
plt.legend()
plt.tight_layout()
plt.savefig("fig_aapl_forecast.png", dpi=120)
plt.show()
\end{lstlisting}

\textbf{Resultados obtenidos:}
\begin{itemize}
  \item $\hat{\phi}_1 = -0.0739$, \quad $\hat{\phi}_2 = -0.8667$
  \item $\hat{\theta}_1 = \phantom{-}0.0568$, \quad $\hat{\theta}_2 = \phantom{-}0.8844$
  \item $\hat{\sigma}^2 = 3.9115$
  \item $\mathrm{AIC} = 5\,489.09$, \quad $\mathrm{BIC} = 5\,514.96$
\end{itemize}

\textbf{Interpretación.} Los coeficientes AR de orden 2 y MA de orden 2 capturan la
dinámica de reversión y memoria corta del precio diario. La suma $\phi_1+\phi_2 \approx -0.94$
indica que la serie diferenciada es casi estacionaria al borde. El $\hat{\sigma}^2 = 3.91$
refleja la volatilidad diaria residual en dólares al cuadrado.

% ─────────────────────────────────────────────────────────
\subsection{Ejemplo 2: Bitcoin (BTC-USD) --- ARIMA$(1,1,1)$ en log-precio}
% ─────────────────────────────────────────────────────────

\textbf{Contexto.} Bitcoin es conocido por su alta volatilidad. Se modela el
\emph{logaritmo} del precio para estabilizar la varianza (transformación Box-Cox con
$\lambda=0$). Los retornos logarítmicos son estacionarios, confirmado por la prueba ADF.

\begin{lstlisting}[caption={Ejemplo 2 --- ARIMA$(1,1,1)$ para BTC en log-precio}]
# ── Datos BTC (log-precio simulado) ─────────────────────
dates_btc = pd.date_range("2021-01-01", "2024-12-31", freq='B')
n2 = len(dates_btc)
btc_log = pd.Series(
    np.cumsum(np.random.randn(n2) * 0.04) + np.log(35000),
    index=dates_btc,
    name="BTC_log"
)

# ── Estacionariedad ──────────────────────────────────────
test_adf(btc_log,            "BTC log-precio") # NO estacionaria
test_adf(btc_log.diff(),     "BTC log-retorno") # estacionaria

# ── ACF y PACF ───────────────────────────────────────────
graficar_acf_pacf(btc_log.diff().dropna(), "BTC log-retorno")

# ── Ajuste ARIMA(1,1,1) ──────────────────────────────────
modelo_btc = SARIMAX(btc_log, order=(1,1,1), trend='n')
res_btc    = modelo_btc.fit(disp=False)
print(f"AIC = {res_btc.aic:.2f}  |  BIC = {res_btc.bic:.2f}")
print(res_btc.params.round(4))

# ── Pronostico 14 dias (escala original USD) ─────────────
fc_btc  = res_btc.get_forecast(steps=14)
ic_btc  = fc_btc.conf_int(alpha=0.05)
precio  = np.exp(fc_btc.predicted_mean.iloc[-1])
p_lo    = np.exp(ic_btc.iloc[-1, 0])
p_hi    = np.exp(ic_btc.iloc[-1, 1])
print(f"BTC dia 14 (USD): ${precio:,.0f}")
print(f"IC 95% (USD): [${p_lo:,.0f},  ${p_hi:,.0f}]")

# ── Grafica retornos log ─────────────────────────────────
plt.figure(figsize=(10,4))
plt.plot(btc_log.diff().dropna().values, linewidth=0.6, color='#ff7f0e')
plt.axhline(0, color='black', linewidth=0.5, linestyle='--')
plt.title("BTC - Log-retornos diarios")
plt.tight_layout()
plt.savefig("fig_btc_returns.png", dpi=120)
plt.show()
\end{lstlisting}

\textbf{Resultados obtenidos:}
\begin{itemize}
  \item $\hat{\phi}_1 = -0.2688$, \quad $\hat{\theta}_1 = 0.2882$
  \item $\hat{\sigma}^2 = 0.0016$ (varianza del log-retorno)
  \item $\mathrm{AIC} = -3\,776.34$, \quad $\mathrm{BIC} = -3\,761.49$
\end{itemize}

\textbf{Interpretación.} El componente AR$(1)$ negativo indica \emph{reversión a la
media} de corto plazo en los retornos: un retorno positivo hoy tiende a ir seguido de
un retorno ligeramente negativo mañana. El MA$(1)$ positivo captura autocorrelación
positiva inmediata. La transformación logarítmica es esencial: sin ella, los residuos
presentarían heterocedasticidad severa violando los supuestos del modelo.

% ─────────────────────────────────────────────────────────
\subsection{Ejemplo 3: S\&P~500 ETF (SPY) --- SARIMA$(1,1,1)(1,1,1)_{12}$}
% ─────────────────────────────────────────────────────────

\textbf{Contexto.} Se utilizan cierres \emph{mensuales} del SPY para detectar
estacionalidad anual. Con $s=12$, el modelo SARIMA captura tanto la tendencia de largo
plazo como los patrones de enero, verano y fin de año documentados en los mercados de
renta variable.

\begin{lstlisting}[caption={Ejemplo 3 --- SARIMA$(1,1,1)(1,1,1)_{12}$ para SPY mensual}]
# ── Datos SPY mensual (log-precio simulado) ──────────────
dates_spy = pd.date_range("2015-01", "2024-12", freq='ME')
n3 = len(dates_spy)
trend3    = np.linspace(np.log(2000), np.log(4800), n3)
seasonal3 = 0.02 * np.sin(2*np.pi*np.arange(n3)/12)
noise3    = np.random.randn(n3) * 0.02
spy_log   = pd.Series(trend3 + seasonal3 + noise3,
                      index=dates_spy, name="SPY_log")

# ── Estacionariedad ──────────────────────────────────────
test_adf(spy_log,        "SPY log-precio") # NO estacionaria
test_adf(spy_log.diff(), "SPY log-retorno") # estacionaria

# ── ACF y PACF (detectar estacionalidad en multiplos de 12) ──
graficar_acf_pacf(spy_log.diff().dropna(), "SPY mensual", lags=24)

# ── Ajuste SARIMA(1,1,1)(1,1,1)[12] ─────────────────────
modelo_spy = SARIMAX(
    spy_log,
    order=(1, 1, 1),
    seasonal_order=(1, 1, 1, 12),
    trend='n'
)
res_spy = modelo_spy.fit(disp=False)
print(f"AIC = {res_spy.aic:.2f}  |  BIC = {res_spy.bic:.2f}")
print(res_spy.params.round(4))

# ── Pronostico 12 meses (escala original USD) ─────────────
fc_spy  = res_spy.get_forecast(steps=12)
ic_spy  = fc_spy.conf_int(alpha=0.05)
precio  = np.exp(fc_spy.predicted_mean.iloc[-1])
p_lo    = np.exp(ic_spy.iloc[-1, 0])
p_hi    = np.exp(ic_spy.iloc[-1, 1])
print(f"SPY en 12 meses (USD): ${precio:,.0f}")
print(f"IC 95% (USD): [${p_lo:,.0f},  ${p_hi:,.0f}]")

# ── Grafica pronostico ───────────────────────────────────
hist_vals = np.exp(spy_log.values)
fc_vals   = np.exp(fc_spy.predicted_mean.values)
ic_lo     = np.exp(ic_spy.iloc[:,0].values)
ic_hi     = np.exp(ic_spy.iloc[:,1].values)

plt.figure(figsize=(10,4))
plt.plot(range(len(hist_vals)), hist_vals, label='Historico')
x_fc = range(len(hist_vals), len(hist_vals)+12)
plt.plot(x_fc, fc_vals, color='red', label='Pronostico')
plt.fill_between(x_fc, ic_lo, ic_hi,
                 alpha=0.2, color='red', label='IC 95%')
plt.title("SPY - SARIMA(1,1,1)(1,1,1)[12]: Pronostico 12 meses")
plt.legend()
plt.tight_layout()
plt.savefig("fig_spy_forecast.png", dpi=120)
plt.show()
\end{lstlisting}

\textbf{Resultados obtenidos:}
\begin{itemize}
  \item AR regular:     $\hat{\phi}_1   = -0.0411$
  \item MA regular:     $\hat{\theta}_1 = -0.9558$
  \item AR estacional:  $\hat{\Phi}_1   = -0.1470$
  \item MA estacional:  $\hat{\Theta}_1 = -0.8795$
  \item $\hat{\sigma}^2 = 0.0004$
  \item $\mathrm{AIC} = -508.29$, \quad $\mathrm{BIC} = -494.97$
\end{itemize}

\textbf{Interpretación.} El MA estacional fuerte ($\hat{\Theta}_1 \approx -0.88$) indica
que el retorno del mes actual está negativamente correlacionado con el del mismo mes del
año anterior, capturando la reversión estacional. Esto es consistente con el
\emph{efecto enero} y la tendencia bajista de septiembre documentada históricamente
en los mercados de renta variable.

% ─────────────────────────────────────────────────────────
\section{Comparación de Modelos}
% ─────────────────────────────────────────────────────────

\begin{table}[H]
\centering
\caption{Resumen comparativo de los tres modelos ajustados}
\begin{tabular}{lllrrr}
\toprule
\textbf{Activo} & \textbf{Modelo} & \textbf{Frecuencia}
  & \textbf{AIC} & \textbf{BIC} & $\hat{\sigma}^2$ \\
\midrule
AAPL & ARIMA$(2,1,2)$                    & Diaria   &  5\,489.09  &  5\,514.96  & 3.9115 \\
BTC  & ARIMA$(1,1,1)$                    & Diaria   & $-$3\,776.34 & $-$3\,761.49 & 0.0016 \\
SPY  & SARIMA$(1,1,1)(1,1,1)_{12}$       & Mensual  & $-$508.29   & $-$494.97   & 0.0004 \\
\bottomrule
\end{tabular}
\end{table}

\noindent\textbf{Nota.} Los valores de AIC y BIC solo son comparables dentro de la
misma serie. No deben usarse para comparar modelos ajustados a distintas series o
distintas frecuencias temporales.

% ─────────────────────────────────────────────────────────
\section{Conclusiones}
% ─────────────────────────────────────────────────────────

\begin{itemize}
  \item El modelo $(S)ARIMA(p,d,q)(P,D,Q)_s$ generaliza de forma elegante los modelos
        AR, MA e ARIMA incluyendo estacionalidad multiplicativa.
  \item La \textbf{ACF} y la \textbf{PACF} son las principales herramientas de
        identificación visual del orden del modelo.
  \item Los criterios \textbf{AIC} y \textbf{BIC} permiten seleccionar el modelo más
        parsimonioso con mejor ajuste.
  \item En series financieras, la \textbf{transformación logarítmica} es prácticamente
        obligatoria para cumplir los supuestos del modelo.
  \item Para volatilidad condicional (varianza no constante), los modelos ARIMA deben
        complementarse con modelos \textbf{GARCH}.
\end{itemize}

% ─────────────────────────────────────────────────────────
\begin{thebibliography}{9}

\bibitem{boxjenkins}
Box, G.E.P., Jenkins, G.M., Reinsel, G.C., \& Ljung, G.M.\ (2015).
\textit{Time Series Analysis: Forecasting and Control} (5th ed.). Wiley.

\bibitem{hamilton}
Hamilton, J.D.\ (1994). \textit{Time Series Analysis}. Princeton University Press.

\bibitem{hyndman}
Hyndman, R.J.\ \& Athanasopoulos, G.\ (2021).
\textit{Forecasting: Principles and Practice} (3rd ed.).
OTexts. \url{https://otexts.com/fpp3/}

\bibitem{statsmodels}
Seabold, S.\ \& Perktold, J.\ (2010).
Statsmodels: Econometric and Statistical Modeling with Python.
\textit{Proceedings of the 9th Python in Science Conference}.

\end{thebibliography}

\end{document}
